\documentclass[11pt,a4paper]{article}
\usepackage{fontspec}
\usepackage{polyglossia}
\setdefaultlanguage{spanish}
\usepackage[margin=2.4cm]{geometry}
\usepackage{graphicx}
\usepackage{xcolor}
\usepackage{booktabs}
\usepackage{titlesec}
\usepackage{enumitem}
\usepackage[colorlinks=true,linkcolor=azul,urlcolor=azul]{hyperref}
\usepackage{tcolorbox}

\definecolor{azul}{HTML}{1F4E79}
\definecolor{gris}{HTML}{4A4A4A}
\definecolor{verde}{HTML}{2E7D32}
\definecolor{rojo}{HTML}{B71C1C}
\definecolor{fondo}{HTML}{F2F6FA}

\titleformat{\section}{\Large\bfseries\color{azul}}{\thesection.}{0.5em}{}
\titleformat{\subsection}{\large\bfseries\color{gris}}{}{0em}{}
\setlist{itemsep=2pt,topsep=4pt}
\linespread{1.08}

\newtcolorbox{caja}[1][]{colback=fondo,colframe=azul,boxrule=0.6pt,arc=2pt,
  left=8pt,right=8pt,top=6pt,bottom=6pt,#1}

\begin{document}

\begin{center}
{\Huge\bfseries\color{azul} Un modelo de lenguaje\\[2pt] que sabe decir «no sé»}\\[10pt]
{\large Informe de resultados --- 18 de agosto de 2026}\\[4pt]
{\color{gris} Maximiliano Speranza --- investigador independiente}\\[2pt]
{\color{gris}\small ORCID 0009-0005-0413-8554}
\end{center}

\vspace{6pt}

\begin{caja}
\textbf{Para quién es este documento.} Está escrito para que se entienda sin saber nada de
inteligencia artificial. Cada término técnico se explica la primera vez que aparece. Los números
que figuran acá son los que ya están medidos y verificados; al final hay una sección con lo que
todavía \emph{no} sabemos, que es igual de importante.
\end{caja}

\section{El problema, en una página}

Cuando un sistema de inteligencia artificial no sabe algo, casi nunca lo dice. Inventa una
respuesta con el mismo tono seguro que usa cuando acierta. A eso se le llama \textbf{alucinación}.

El daño no es la respuesta equivocada en sí. Es lo que esa respuesta le hace a todas las demás:

\begin{caja}
Un error silencioso no cuesta una respuesta. Cuesta \textbf{la confianza en todas las otras}.
Si el sistema inventa una de cada diez veces y nunca avisa cuál, hay que verificar las diez.
\end{caja}

Este trabajo ataca tres cosas que en realidad son la misma:

\begin{enumerate}[label=\textbf{\arabic*.}]
  \item Que el modelo \textbf{no olvide} lo que se le dijo antes.
  \item Que \textbf{sepa decir «no sé»} cuando no lo tiene.
  \item Que \textbf{no invente} cuando no sabe.
\end{enumerate}

Son la misma porque un modelo que recuerda bien, pero no distingue entre «me acuerdo» y «no me
acuerdo», va a inventar exactamente en los casos en que olvidó.

\section{Qué se construyó}

Un modelo de lenguaje entrenado \textbf{desde cero}, deliberadamente diminuto:

\begin{center}
\begin{tabular}{ll}
\toprule
Tamaño & 863.730 parámetros --- 3,5 MB \\
Vocabulario & 242 palabras \\
Idioma & artificial, cerrado, pero legible \\
Entrenamiento & desde cero, sin partir de ningún modelo existente \\
\bottomrule
\end{tabular}
\end{center}

Para comparar: los modelos comerciales tienen entre cien mil y un millón de veces más
parámetros. Un \textbf{parámetro} es un número que el modelo ajusta mientras aprende; son las
perillas internas que definen su comportamiento.

\textbf{¿Por qué tan chico a propósito?} Porque acá no se busca que sepa mucho, sino poder
\emph{ver} qué hace por dentro. En un modelo gigante, averiguar por qué inventó algo es
prácticamente imposible: hay demasiadas piezas. En uno de 3,5 MB se puede abrir, medir cada
parte y hacer experimentos controlados en horas en vez de meses.

\subsection{La diferencia con lo que se hace habitualmente}

La forma normal de que un sistema «recuerde» es pegarle un buscador por afuera: cuando llega una
pregunta, el buscador va a una base de datos, trae textos parecidos y se los pasa al modelo. El
modelo, en sí mismo, no recuerda nada.

Acá el archivo de memoria está \textbf{adentro de la arquitectura}, y se entrena junto con el
resto. El modelo no recibe textos: escribe sus propios recuerdos y aprende a consultarlos. Eso
le permite contestar sobre algo dicho en una conversación anterior \textbf{cuyo estado interno ya
se borró}, y si el dato se corrigió después, contestar con la versión corregida.

\section{Los cinco resultados que ya están asegurados}

\subsection{1. El modelo nunca olvida escribir: se equivoca al leer}

Cuando el modelo se equivoca de dato, ¿es porque no lo guardó, o porque no lo encontró?

Se revisaron 8000 casos. \textbf{El dato buscado estaba guardado en el archivo en el 100\,\% de
ellos} --- literalmente ni un caso de 8000 en que faltara.

Cuando responde mal, el dato correcto está ahí, en el segundo puesto de su propia lista de
candidatos, y sólo gana entre el 14 y el 18\,\% de las veces. Cuando responde bien, está en el
primer puesto.

\begin{caja}
\textbf{Por qué esto es una buena noticia:} un error de escritura sería irreparable --- lo que no
se guardó no se puede recuperar. Un error de \emph{lectura} es distinto: la información está,
sólo falta reconocer que la búsqueda salió mal. \textbf{Y eso se puede convertir en un «no sé».}
\end{caja}

\subsection{2. Cuando dice «no sé», es porque fue a buscar y no encontró}

Un modelo podría aprender a decir «no sé» por el motivo equivocado: reconociendo la \emph{forma}
de la pregunta en lugar de consultar su memoria. Sería un truco, no memoria.

El control: se toma la misma pregunta, en el mismo contexto, y se cambia una sola cosa --- si el
dato está guardado o no.

\begin{center}
\begin{tabular}{lr}
\toprule
\textbf{Situación} & \textbf{Veces que dice «no sé»} \\
\midrule
El dato no está en su memoria & $\approx$ 90\,\% \\
El dato sí está & \textbf{0,6\,\% a 2,2\,\%} \\
\bottomrule
\end{tabular}
\end{center}

Con el dato disponible deja de abstenerse y responde, acertando el valor correcto entre el 98 y
el 99\,\% de las veces. \textbf{Está consultando su memoria de verdad.}

\subsection{3. El hallazgo principal: la abstención necesita su propia salida}

Acá está el resultado central, y es de \emph{arquitectura} --- de cómo está construido el modelo,
no de cuánto se lo entrena.

Originalmente, «no sé» era una palabra más del vocabulario. Para responder, el modelo elegía
entre todas sus palabras: «ana», «beto», «743»\ldots\ y «no sé». Todas competían en la misma
votación interna.

El problema es que ahí se mezclan dos preguntas de naturaleza distinta:

\begin{itemize}
  \item «¿Tengo el dato?» --- una pregunta de sí o no, más o menos balanceada.
  \item «¿Cuál es el valor?» --- una elección entre cien opciones.
\end{itemize}

La propuesta fue \textbf{separarlas}: darle a la abstención una salida propia, independiente de
la votación por el valor. Cuesta 129 parámetros más sobre 863.730 --- un 0,015\,\%.

Se probaron tres versiones sobre los mismos modelos, con el mismo entrenamiento y el mismo
presupuesto, cambiando \emph{únicamente} dónde se toma la decisión de callarse:

\begin{center}
\begin{tabular}{lcc}
\toprule
\textbf{Versión} & \textbf{Modelos que pasan la prueba} & \textbf{Error medio} \\
\midrule
\textcolor{verde}{\textbf{Salida propia}} & \textcolor{verde}{\textbf{4 de 5}} & \textcolor{verde}{\textbf{0,093}} \\
Palabra más del vocabulario & 0 de 5 & 0,186 \\
Palabra, pero «amplificada» & 0 de 5 & 0,239 \\
\bottomrule
\end{tabular}
\end{center}

\textbf{De 0 de 5 a 4 de 5 cambiando dónde se decide, no cuánto se entrena.}

El «error medio» de la tabla es la \textbf{falsa abstención}: cuántas veces dice «no sé» aunque
sí tenía el dato. Es el precio de enseñarle a callarse, y es el que hay que mantener bajo ---
porque un modelo que contesta «no sé» a todo nunca se equivoca, y tampoco sirve para nada.

\subsection{4. No era el tamaño: era la competencia}

La tercera versión de la tabla merece una explicación, porque su fracaso enseña algo.

Se había observado que el vector interno de «no sé» era tres veces más corto que el de una
palabra normal --- como si compitiera en desventaja. La versión «amplificada» lo agrandaba para
emparejarlo.

Fracasó en los cinco modelos, y peor que no hacer nada. Al ir a mirar por qué, apareció que la
observación original mezclaba dos partes distintas del modelo: en la parte que \emph{de verdad}
decide la respuesta, «no sé» no estaba en desventaja --- era casi el \textbf{doble} de largo que
una palabra normal.

\begin{caja}
\textbf{Conclusión:} el problema nunca fue el tamaño del vector. Era que dos decisiones distintas
compartían la misma votación. Emparejar los tamaños no cambia nada; separar las decisiones sí.
\end{caja}

\subsection{5. Lo que falta es calibración, no capacidad}

De los cinco modelos, uno no llegó a pasar la prueba. ¿Le falta capacidad, o sólo está mal
ajustado?

El modelo produce internamente un número que indica «cuán seguro estoy de que no lo tengo», y
después lo compara con un punto de corte para decidir si se calla. Ese punto de corte estaba
fijado en el valor de manual, sin ajustar.

Midiendo qué tan bien ese número separa los casos en que el dato está de los casos en que no
está, la puntuación va de \textbf{0,77 a 0,99} (donde 0,5 sería no distinguir nada y 1,0 sería
perfecto). \textbf{La señal está y es buena.} Ajustando el punto de corte con datos separados de
los de la evaluación, pasan 5 de 6 modelos --- incluido el que fallaba.

\section{Cómo sabemos que esto es cierto}

\subsection{Todo se predice por escrito antes de medir}

Cada experimento se escribe primero: qué se espera, con qué números exactos se considera
confirmado y con cuáles refutado. Ese documento se sella con una \textbf{huella digital
criptográfica} --- un código que cambia si se altera una sola letra --- y se publica con fecha
antes de correr nada.

Sirve para una cosa concreta: hace imposible acomodar el criterio después de ver los resultados.
Es una práctica común en medicina y rara en inteligencia artificial.

\subsection{Se reportan también los resultados incómodos}

En este mismo informe hay dos ejemplos:

\begin{itemize}
  \item Una de las cuatro predicciones \textbf{no se cumplió} en su forma literal. El desvío fue
  \emph{a favor} del resultado --- el modelo con salida propia recuerda mejor de lo permitido por
  el criterio --- pero queda registrado como incumplimiento, sin cambiar el criterio después.
  \item La primera versión de un análisis dio un resultado demasiado bueno porque el ajuste y la
  medición usaban los mismos datos. Se rehizo separándolos y el resultado bajó.
\end{itemize}

A lo largo del proyecto se detectaron \textbf{ocho} casos en que un número limpio escondía un
error de instrumentación. Todos están documentados. La costumbre de buscarlos es parte del
método.

\section{Lo que todavía no sabemos}

\begin{itemize}
  \item \textbf{Es un idioma artificial de 242 palabras}, no lenguaje humano. Que el mecanismo
  funcione acá no garantiza que funcione a escala real: lo hace plausible y medible.
  \item \textbf{Un modelo de cinco no pasa} con el ajuste de fábrica. Hay evidencia de que es
  calibración, pero se midió después de ver los resultados, así que es una pista fuerte, no una
  conclusión.
  \item \textbf{No sabemos cuán entrenado tiene que estar} el modelo antes de poder enseñarle a
  callarse. Si hiciera falta entrenarlo casi hasta la perfección, el método no serviría en la
  práctica. Ese experimento está corriendo mientras se escribe este informe.
  \item \textbf{El ajuste del punto de corte no está probado} como parte de un experimento
  predicho de antemano. Requiere su propio registro previo.
\end{itemize}

\section{Por qué importa}

Hoy la industria intenta que los modelos no inventen \textbf{pidiéndoselo}: instrucciones,
filtros, verificadores externos. Todo por afuera del modelo.

Lo que estos resultados sugieren es distinto: \textbf{callarse a tiempo puede ser una propiedad
de cómo está construido el modelo}. No hace falta que sea más grande ni que se lo entrene más.
Hace falta que la decisión de «no sé» no compita con la decisión de «qué respondo».

Ciento veintinueve parámetros sobre ochocientos sesenta mil. Un 0,015\,\% más de modelo, y pasa
de no poder callarse nunca a poder hacerlo casi siempre.

\vfill
\begin{center}
{\color{gris}\small
Todos los experimentos citados están registrados con huella criptográfica y fecha anterior a su
ejecución.\\
Documento generado el 18 de agosto de 2026.}
\end{center}

\end{document}
